% ============================================================
% Capítulo 7: Modelos Predictivos de Trayectoria Académica
% ============================================================
\chapter{Modelos Predictivos de Trayectoria Académica}
\label{ch:modelos_predictivos}

El presente capítulo constituye la etapa final del análisis cuantitativo de esta investigación.
Habiendo establecido en el Capítulo~\ref{ch:supervivencia} la dinámica temporal del egreso mediante
curvas de supervivencia y el modelo de Cox, se implementa aquí un enfoque de aprendizaje automático
supervisado para predecir el \emph{estado académico final} de cada estudiante, formalizando así el
concepto de sistema de alerta estudiantil como herramienta institucional prospectiva.

\section{Marco Conceptual y Diseño Experimental}

Se entrenaron dos modelos de clasificación multiclase basados en el algoritmo
\emph{Random Forest} \citep[cf.][]{breiman2001random}, cada uno diferenciado por el horizonte
temporal de sus predictores:

\begin{itemize}
    \item \textbf{Modelo~A — Alerta Temprana (S1–S4):} utiliza exclusivamente las métricas de
          regularidad, esfuerzo, promedio natural y promedio real de los primeros cuatro semestres,
          junto con la variable de interacción $\text{esf}_i \times \text{real}_i$. Simula un
          escenario institucional en el que se desea identificar estudiantes en riesgo a la mitad
          de la trayectoria curricular.
    \item \textbf{Modelo~B — Trayectoria Completa (S1–S8):} extiende el conjunto de predictores
          hasta el octavo semestre. Representa el horizonte temporal máximo en que la información
          académica aún resulta accionable para intervención institucional.
\end{itemize}

La variable objetivo $Y$ toma cuatro valores categóricos:
\emph{Titulado a tiempo}, \emph{Titulado con rezago}, \emph{Abandono silencioso} y
\emph{Activo/Aún cursando}. La clasificación se realizó sobre una muestra combinada de ambas
carreras ($N_A = 4{,}853$; $N_B = 4{,}276$ observaciones completas respectivamente).

\subsection{Protocolo de Validación}

Para evitar el sobreajuste en la evaluación, se utilizó \textbf{validación cruzada estratificada
de 5 pliegues} (\emph{Stratified 5-fold CV}), generando predicciones \emph{out-of-fold} (OOF)
sobre el conjunto completo. Los hiperparámetros del modelo se fijaron en:
$n_{\text{árboles}} = 500$, $\text{max\_features} = \sqrt{p}$, $\text{min\_samples\_leaf} = 3$,
ponderación de clases balanceada (\texttt{class\_weight=`balanced'}), semilla aleatoria $= 42$.

Los \textbf{intervalos de confianza al 95\,\%} para todas las métricas de desempeño se obtuvieron
mediante \textbf{bootstrapping no paramétrico} con $B = 2{,}000$ remuestreos sobre las
predicciones OOF, empleando percentiles empíricos al 2.5\,\% y 97.5\,\%.

% ──────────────────────────────────────────────────────────────────────────────
\section{Resultados de Desempeño Global}
% ──────────────────────────────────────────────────────────────────────────────

La Tabla~\ref{tab:rf_global} sintetiza las métricas de desempeño global con sus respectivos
intervalos de confianza. El Modelo~B supera consistentemente al Modelo~A en todas las métricas.

\begin{table}[H]
\centering
\caption{Desempeño global de los modelos Random Forest (Validación cruzada OOF, $B=2{,}000$ bootstrap)}
\label{tab:rf_global}
\begin{tabular}{lcccc}
\toprule
\textbf{Modelo} & \textbf{$N$} & \textbf{Accuracy} & \textbf{IC 95\,\% Accuracy} & \textbf{F1-macro} \\
\midrule
A — Alerta Temprana (S1–S4) & 4,853 & 0.655 & [0.641, 0.669] & 0.513 \\
B — Trayectoria Completa (S1–S8) & 4,276 & 0.782 & [0.770, 0.794] & 0.591 \\
\bottomrule
\end{tabular}
\end{table}

\noindent
El Modelo~A alcanza una exactitud de $0.655$ (IC~95\,\%: $[0.641, 0.669]$) usando únicamente
información de los primeros cuatro semestres. El Modelo~B mejora esta cifra a $0.782$
(IC~95\,\%: $[0.770, 0.794]$) al incorporar la trayectoria hasta el octavo semestre.
En ambos casos los intervalos no se solapan, indicando diferencias estadísticamente robustas.

% ──────────────────────────────────────────────────────────────────────────────
\section{Desempeño por Clase}
% ──────────────────────────────────────────────────────────────────────────────

La Tabla~\ref{tab:rf_clases} desglosa la precisión y el recall por categoría académica, con
sus intervalos de confianza bootstrap al 95\,\%.

\begin{table}[H]
\centering
\caption{Precisión y Recall por clase — IC al 95\,\% (Bootstrap, $B=2{,}000$)}
\label{tab:rf_clases}
\scriptsize
\setlength{\tabcolsep}{4pt}
\begin{tabular}{llcccc}
\toprule
\textbf{Modelo} & \textbf{Clase} & \textbf{Precisión} & \textbf{IC 95\,\% Precisión} & \textbf{Recall} & \textbf{IC 95\,\% Recall} \\
\midrule
\multirow{4}{*}{A (S1–S4)}
 & Abandono silencioso    & 0.587 & [0.560, 0.615] & 0.651 & [0.624, 0.680] \\
 & Activo / Aún cursando & 0.204 & [0.098, 0.324] & 0.039 & [0.017, 0.064] \\
 & Titulado a tiempo     & 0.785 & [0.767, 0.803] & 0.800 & [0.782, 0.818] \\
 & Titulado con rezago   & 0.568 & [0.544, 0.593] & 0.585 & [0.559, 0.610] \\
\midrule
\multirow{4}{*}{B (S1–S8)}
 & Abandono silencioso    & 0.623 & [0.591, 0.654] & 0.712 & [0.681, 0.743] \\
 & Activo / Aún cursando & 0.169 & [0.031, 0.333] & 0.022 & [0.005, 0.045] \\
 & Titulado a tiempo     & 0.901 & [0.886, 0.914] & 0.904 & [0.890, 0.918] \\
 & Titulado con rezago   & 0.747 & [0.726, 0.768] & 0.769 & [0.748, 0.788] \\
\bottomrule
\end{tabular}
\end{table}

\noindent
En ambos modelos, la categoría \emph{Titulado a tiempo} presenta el mejor desempeño, con una
precisión que alcanza $0.901$ (IC: $[0.886, 0.914]$) en el Modelo~B. En contraste, la clase
\emph{Activo/Aún cursando} resulta la más difícil de identificar en ambos escenarios (recall
$< 0.04$), lo que es esperable dado su carácter ambiguo: los estudiantes activos aún no han
expresado su estado final real.

% ──────────────────────────────────────────────────────────────────────────────
\section{Matrices de Confusión}
% ──────────────────────────────────────────────────────────────────────────────

Las Figuras~\ref{fig:cm_s4} y~\ref{fig:cm_s8} presentan las matrices de confusión OOF de
ambos modelos.

\begin{figure}[H]
    \centering
    \begin{subfigure}[b]{0.48\textwidth}
        \centering
        \includegraphics[width=\textwidth]{rf_confusion_matrix_s4}
        \caption{Modelo A — Alerta Temprana (S1–S4)}
        \label{fig:cm_s4}
    \end{subfigure}
    \hfill
    \begin{subfigure}[b]{0.48\textwidth}
        \centering
        \includegraphics[width=\textwidth]{rf_confusion_matrix_s8}
        \caption{Modelo B — Trayectoria Completa (S1–S8)}
        \label{fig:cm_s8}
    \end{subfigure}
    \caption{Matrices de confusión (validación cruzada OOF, 5-fold estratificado).
             Diagonal principal: predicciones correctas.}
    \label{fig:confusion_matrices}
\end{figure}

\noindent
La principal fuente de error en ambos modelos es la confusión entre \emph{Abandono silencioso}
y \emph{Titulado con rezago}: en el Modelo~A, 324 estudiantes en situación de abandono son
clasificados erróneamente como rezago, y 302 en situación de rezago son clasificados como
abandono. El Modelo~B reduce estas confusiones ($195$ y $194$ respectivamente), confirmando
que los semestres avanzados son informativos para discriminar ambas categorías.

% ──────────────────────────────────────────────────────────────────────────────
\section{Importancia de Variables}
% ──────────────────────────────────────────────────────────────────────────────

Las Figuras~\ref{fig:fi_s4} y~\ref{fig:fi_s8} muestran los doce predictores con mayor
importancia Gini en cada modelo.

\begin{figure}[H]
    \centering
    \begin{subfigure}[b]{0.48\textwidth}
        \centering
        \includegraphics[width=\textwidth]{rf_feature_importance_s4}
        \caption{Modelo A — S1–S4}
        \label{fig:fi_s4}
    \end{subfigure}
    \hfill
    \begin{subfigure}[b]{0.48\textwidth}
        \centering
        \includegraphics[width=\textwidth]{rf_feature_importance_s8}
        \caption{Modelo B — S1–S8}
        \label{fig:fi_s8}
    \end{subfigure}
    \caption{Importancia de variables (Gini) para los modelos Random Forest.
             Código de colores: \textcolor{red}{rojo} = interacción esf$\times$real,
             \textcolor{blue}{azul} = esfuerzo, \textcolor{green!60!black}{verde} = regularidad,
             \textcolor{orange}{naranja} = promedio.}
    \label{fig:feature_importances}
\end{figure}

\noindent
En el \textbf{Modelo~A}, la regularidad del cuarto semestre (\texttt{s4\_reg}, importancia
$= 0.126$) emerge como el predictor más influyente, seguido del promedio real del cuarto
semestre (\texttt{s4\_real}, $= 0.095$) y la interacción $\text{esf}_4 \times \text{real}_4$
($= 0.084$). La regularidad del tercer semestre (\texttt{s3\_reg}, $= 0.081$) también tiene
peso relevante, sugiriendo que la inercia académica de los primeros semestres empieza a
manifestarse a partir de S3.

En el \textbf{Modelo~B}, la regularidad del octavo semestre (\texttt{s8\_reg}, $= 0.130$)
domina el árbol de decisión, seguida de la regularidad del séptimo (\texttt{s7\_reg}, $= 0.089$)
y la interacción $\text{esf}_8 \times \text{real}_8$ ($= 0.065$). Este patrón confirma que
la regularidad acumulada es el indicador más robusto del destino académico, consistente con
los hallazgos del modelo de Cox del capítulo anterior.

% ──────────────────────────────────────────────────────────────────────────────
\section{Prueba de Hipótesis: Test de McNemar}
% ──────────────────────────────────────────────────────────────────────────────

Para establecer con rigor estadístico si la diferencia en desempeño entre ambos modelos
es significativa o atribuible a variación aleatoria, se aplicó la \textbf{prueba de McNemar}
\citep{mcnemar1947note} sobre los pares de predicciones comunes a ambos modelos.

\vspace{0.5em}
\noindent\textbf{Hipótesis:}
\begin{itemize}
    \item[$H_0$:] Los modelos A y B tienen la misma tasa de error (las discordancias son simétricas).
    \item[$H_1$:] Los modelos difieren significativamente en su capacidad predictiva.
\end{itemize}

La prueba opera sobre la tabla de contingencia $2\times 2$ de discordancias:
$b = 1{,}815$ (A correcto, B incorrecto) y $c = 446$ (A incorrecto, B correcto).
La asimetría marcada ($b \gg c$) indica que el Modelo~B corrige sistemáticamente errores
que el Modelo~A no puede resolver. Se aplicó la corrección de continuidad de Yates.

\begin{table}[H]
\centering
\caption{Resultado de la Prueba de McNemar (Modelo~A vs Modelo~B)}
\label{tab:mcnemar}
\begin{tabular}{lcc}
\toprule
\textbf{Parámetro} & \textbf{Valor} \\
\midrule
$b$ (A correcto, B incorrecto) & 1,815 \\
$c$ (A incorrecto, B correcto) & 446 \\
Estadístico $\chi^2$ (Yates)   & 827.70 \\
$p$-valor                       & $5.13 \times 10^{-182}$ \\
\midrule
\textbf{Decisión} ($\alpha = 0.05$) & \textbf{Rechazar $H_0$} \\
\bottomrule
\end{tabular}
\end{table}

\noindent
El valor del estadístico $\chi^2 = 827.70$ con un $p\text{-valor} = 5.13 \times 10^{-182}$
proporciona evidencia abrumadora para rechazar $H_0$. La incorporación de los semestres S5–S8
no es trivial: produce una mejora en predicción estadísticamente significativa y cuantitativamente
sustantiva ($+12.7$ puntos porcentuales en accuracy).

% ──────────────────────────────────────────────────────────────────────────────
\section{Síntesis e Implicaciones Institucionales}
% ──────────────────────────────────────────────────────────────────────────────

\begin{figure}[H]
    \centering
    \includegraphics[width=0.85\textwidth]{rf_comparativa_ic}
    \caption{Comparativa de Accuracy y F1-Score macro con intervalos de confianza al 95\,\%
             (Bootstrap). Los intervalos no se solapan en ninguna métrica.}
    \label{fig:comparativa_ic}
\end{figure}

Los resultados permiten derivar tres conclusiones de relevancia institucional:

\begin{enumerate}
    \item \textbf{Valor del sistema de alerta temprana.} Con solo cuatro semestres de información,
          el Modelo~A clasifica correctamente al $65.5\,\%$ de los estudiantes, con una precisión
          del $80\,\%$ para el grupo más crítico (\emph{Titulado a tiempo}). Esto justifica la
          implementación de un panel de riesgo temprano en los sistemas de gestión escolar.

    \item \textbf{Ganancia marginal de la información tardía.} El Modelo~B mejora la clasificación
          en $+12.7$ puntos porcentuales de accuracy y $+7.8$ puntos de F1-macro. La prueba de
          McNemar confirma que esta ganancia es estadísticamente significativa
          ($p = 5.13\times10^{-182}$), no un artefacto muestral.

    \item \textbf{Límites del modelo: la clase ``Activo/Aún cursando''.} Ambos modelos tienen
          dificultad para identificar estudiantes activos (recall $< 5\,\%$). Esto es inherente
          al problema: los estudiantes activos no han expresado su trayectoria completa, por lo que
          su patrón de comportamiento se solapa con el de estudiantes que eventualmente se titularán
          o abandonarán. Una solución prospectiva requeriría datos longitudinales actualizados
          semestre a semestre.
\end{enumerate}

\noindent
Finalmente, la consistencia de la \textbf{regularidad semestral} como el predictor más influyente
en ambos modelos —confirmando los hallazgos del análisis de supervivencia— refuerza la hipótesis
central de esta investigación: la \emph{inercia de inscripción continua} es el mecanismo
estadístico dominante en la determinación del destino académico de los estudiantes de Ciencias.

% ══════════════════════════════════════════════════════════════════════════════
\section{Enfoque Truncado: Regresión Regularizada sin Censura}
\label{sec:enfoque_truncado}
% ══════════════════════════════════════════════════════════════════════════════

Mientras que el Enfoque Multiclase trata el estado académico como variable categórica,
el \textbf{Enfoque Truncado} aborda el problema como una regresión continua. Al restringir
el análisis a las generaciones anteriores a 2017 ($\leq 2016$) y eliminar registros sin
semestre de término conocido, se obtiene una muestra \emph{libre de censura} que permite
aplicar regresión regularizada sin los sesgos de tratar valores censurados como numéricos.

\subsection{Datos y Protocolo}

La muestra comprende $N = 3{,}157$ estudiantes ($\bar{y} = 11.29$ semestres, $\sigma = 2.75$;
rango: $[8, 20]$), con $p = 40$ predictores: las cuatro métricas del Vector Excalibur
(regularidad, esfuerzo, promedio natural, promedio real) en los semestres S1–S8, más las
interacciones $\text{esf}_i \times \text{real}_i$. Los predictores se estandarizaron (Z-score)
antes del entrenamiento. Se entrenaron dos modelos con $\alpha$ óptimo por 5-fold CV:
\textbf{Ridge} ($\alpha^* = 7.743$, 40 coefs. activos) y \textbf{Lasso}
($\alpha^* = 0.0024$, 31 coefs. activos). Los IC al 95\,\% se obtuvieron por bootstrapping
($B = 2{,}000$).

\subsection{Resultados de Desempeño}

\begin{table}[H]
\centering
\caption{Desempeño del Enfoque Truncado — Ridge vs.\ Lasso (IC al 95\,\%, Bootstrap $B=2{,}000$)}
\label{tab:trunc_global}
\begin{tabular}{lcccc}
\toprule
\textbf{Modelo} & \textbf{RMSE (IC 95\,\%)} & \textbf{MAE (IC 95\,\%)} & \textbf{R² (IC 95\,\%)} & \textbf{Activos} \\
\midrule
Ridge ($\alpha = 7.743$)  & 1.360 [1.307, 1.421] & 0.957 [0.924, 0.991] & 0.755 [0.738, 0.772] & 40/40 \\
Lasso ($\alpha = 0.0024$) & 1.360 [1.307, 1.421] & 0.957 [0.925, 0.992] & 0.755 [0.738, 0.772] & 31/40 \\
\bottomrule
\end{tabular}
\end{table}

\noindent
Ambos modelos alcanzan $R^2 \approx 0.755$, explicando el $75.5\,\%$ de la varianza con un
Error Absoluto Medio de $\approx 0.957$ semestres —es decir, el modelo predice el semestre
exacto de egreso con un error inferior a \textbf{un semestre} en promedio.

\begin{figure}[H]
    \centering
    \includegraphics[width=\textwidth]{trunc_pred_vs_real}
    \caption{Predicción vs.\ Real del semestre de término (Ridge izquierda, Lasso derecha).
             Los puntos sobre la diagonal indican predicciones perfectas.}
    \label{fig:trunc_pred_vs_real}
\end{figure}

\subsection{Prueba F Global y Coeficientes Significativos}

La \textbf{prueba F global} rechaza $H_0$ con certeza práctica para ambos modelos:
$F(40,\, 3116) = 240.05$, $p < 2.2 \times 10^{-16}$.

Los coeficientes individualmente significativos ($p < 0.05$) bajo Ridge se presentan
en la Tabla~\ref{tab:trunc_coef_sig}.

\begin{table}[H]
\centering
\caption{Coeficientes significativos del modelo Ridge (prueba $t$, $\alpha = 0.05$)}
\label{tab:trunc_coef_sig}
\begin{tabular}{lrrrrl}
\toprule
\textbf{Variable} & \textbf{Coef.} & \textbf{SE} & \textbf{$t$} & \textbf{$p$-valor} & \\
\midrule
\texttt{s8\_reg}          & $-2.445$ & 0.179 & $-13.64$ & $< 10^{-40}$ & ✓ \\
\texttt{s7\_reg}          & $-0.818$ & 0.276 & $-2.965$ & $0.003$ & ✓ \\
\texttt{s5\_esf}          & $+0.329$ & 0.130 & $+2.525$ & $0.012$ & ✓ \\
\texttt{s7\_esf}          & $+0.289$ & 0.130 & $+2.216$ & $0.027$ & ✓ \\
\texttt{s6\_nat}          & $+0.339$ & 0.158 & $+2.141$ & $0.032$ & ✓ \\
\texttt{s5\_esf\_x\_real} & $-0.377$ & 0.188 & $-2.001$ & $0.045$ & ✓ \\
\bottomrule
\end{tabular}
\end{table}

\noindent
El coeficiente de \texttt{s8\_reg} $= -2.445$ ($t = -13.64$) es el más significativo del
modelo: cada unidad adicional de regularidad en el octavo semestre \emph{reduce} el semestre
de término predicho en $2.4$ semestres, controlando por el resto. Esta es la cuantificación
más precisa del mecanismo de inercia académica documentado previamente.

\begin{figure}[H]
    \centering
    \begin{subfigure}[b]{0.48\textwidth}
        \centering
        \includegraphics[width=\textwidth]{trunc_ridge_importancia}
        \caption{Ridge — Top 12 predictores ($|\beta|$ estandarizado)}
    \end{subfigure}
    \hfill
    \begin{subfigure}[b]{0.48\textwidth}
        \centering
        \includegraphics[width=\textwidth]{trunc_lasso_importancia}
        \caption{Lasso — Top 12 predictores (solo activos)}
    \end{subfigure}
    \caption{Importancia de variables — Enfoque Truncado. \texttt{s8\_reg} domina en ambos modelos,
             seguido de \texttt{s7\_reg} y las métricas del quinto semestre.}
    \label{fig:trunc_importancias}
\end{figure}

\subsection{Prueba Diebold-Mariano: Ridge vs.\ Lasso}

La \textbf{prueba de Diebold-Mariano} \citep{diebold1995comparing} no rechaza $H_0$ de
igualdad de errores ($DM = 0.026$, $p = 0.979$). La parsimonia de Lasso —que elimina
9 predictores redundantes— no sacrifica desempeño estadísticamente significativo.
Por interpretabilidad institucional, se prefiere el \textbf{modelo Lasso}.

\subsection{Convergencia con el Enfoque Multiclase}

La coherencia entre ambos enfoques consolida las conclusiones del capítulo:

\begin{enumerate}
    \item \textbf{Regularidad como predictor dominante.} En Random Forest la importancia Gini
          de \texttt{s8\_reg} es la más alta; en Lasso/Ridge su coeficiente tiene el mayor
          $|t|$ por un margen de $5\times$.
    \item \textbf{Semestres tardíos son más informativos.} La ganancia del Modelo~B sobre A
          y la dominancia de \texttt{s7\_reg}/\texttt{s8\_reg} en regresión apuntan al mismo
          fenómeno: la trayectoria \emph{cristaliza} en S7–S8.
    \item \textbf{El promedio individual tiene efecto marginal.} Ningún predictor de promedio
          (nat/real) es significativo en la regresión, confirmando cuantitativamente la paradoja
          del perfeccionismo documentada en el Capítulo~\ref{ch:supervivencia}.
\end{enumerate}
